\documentclass[11pt]{article}

\usepackage{times}
\usepackage[margin=1in]{geometry}
\usepackage{graphicx}
\usepackage{booktabs}
\usepackage{amsmath}
\usepackage{hyperref}
\usepackage{caption}
\usepackage{subcaption}
\usepackage{float}
\usepackage{parskip}
\usepackage{setspace}

\hypersetup{
    colorlinks=true,
    linkcolor=blue,
    urlcolor=blue
}

\title{\textbf{Network Flow Forensics with Machine Learning}\\
       \large CSCI 4540 --- Final Project Option \#2}
\author{Jhelani Stokes}
\date{April 27, 2026}

\begin{document}

\maketitle
\thispagestyle{empty}

% ────────────────────────────────────────────────────────────────────────────
\section{Introduction}

This project applies supervised machine learning
to a dataset of 2,500 network flow records, labelled as either \emph{benign} or
\emph{malicious}, with the goal of building accurate classifiers and critically
evaluating their forensic utility.

This report is structured as follows: Section~\ref{sec:eda}
describes the dataset and exploratory analysis; Section~\ref{sec:features}
details feature engineering; Section~\ref{sec:models} covers model training;
Section~\ref{sec:eval} presents evaluation results and cross-validation
comparison; and Section~\ref{sec:forensics} discusses forensic implications and
ethical concerns. All code is available at the GitHub repository linked in
Section~\ref{sec:code}.

% ────────────────────────────────────────────────────────────────────────────
\section{Data Exploration and Preprocessing}
\label{sec:eda}

\subsection{Dataset Overview}

The dataset contains 2,500 flow records with ten original features and one
binary target variable (\texttt{is\_malicious}). Features cover four
categories: temporal (\texttt{duration}), volumetric (\texttt{src\_bytes},
\texttt{dst\_bytes}, \texttt{packet\_count}, \texttt{avg\_packet\_size},
\texttt{flow\_rate}), reliability (\texttt{num\_failed\_connections}), and
protocol type (one-hot encoded flags for TCP, UDP, and ICMP).

No missing values were detected across any feature. Descriptive statistics
reveal substantial right-skew in volumetric fields: \texttt{src\_bytes} has
a mean of 2,192 bytes but a maximum of 78,420 bytes, indicating the presence
of high-volume outlier flows. 

\subsection{Class Distribution}

Of the 2,500 records, 1,738 (69.5\%) are labelled benign and 762 (30.5\%)
malicious. To account for this imbalance, all models were trained with \texttt{class\_weight="balanced"}, which more heavily penalizes the model for misclassifying rarer classes in the dataset.

\subsection{Feature Correlation}

A correlation heatmap (Appendix, Figure~\ref{fig:heatmap}) created with \texttt{pd.corr} highlights several interesting relationships between the input features and classification. The feature \texttt{packet\_count} shows the highest correlation with malicious traffic. This is not very surprising because a high packet count is observed when a network is being scanned. Between the three protocols ICMP had almost no correlation, TCP showed negative correlation meaning that TCP traffic was indicative of non-malicious activity, and UDP was moderately correlated with malicious activity.
% ────────────────────────────────────────────────────────────────────────────
\section{Feature Engineering}
\label{sec:features}

3 engineered features were added to the original ten:

\begin{itemize}
    \item \textbf{byte\_ratio} $= \text{src\_bytes}\,/\,(\text{dst\_bytes}+1)$:
          Asymmetric upload/download ratios are potential indicators of malicious traffic.
    \item \textbf{bytes\_per\_packet} $= (\text{src\_bytes}+\text{dst\_bytes})\,/\,(\text{packet\_count}+1)$:
          Unusually small or large values can be a sign of network scans or anomalous bulk transfers.
    \item \textbf{failure\_rate} $= \text{num\_failed\_connections}\,/\,(\text{packet\_count}+1)$:
          Normalises raw failure counts relative to flow size.
\end{itemize}

After engineering, the feature matrix has 13 columns. All continuous features
were also standardized inside a
\texttt{scikit-learn} \texttt{Pipeline}.

% ────────────────────────────────────────────────────────────────────────────
\section{Model Training}
\label{sec:models}

Two classifiers were implemented, each wrapped in a \texttt{Pipeline}
with a \texttt{StandardScaler} prefix:

\begin{enumerate}
    \item \textbf{Logistic Regression (LR)}: A linear baseline well-suited to
          linearly separable high-dimensional problems.
    \item \textbf{Random Forest (RF)}: An ensemble of 200 decision trees that
          captures non-linear feature interactions and provides interpretable
          feature importances.
\end{enumerate}

The dataset was split 80/20 into training (2,000 records) and test
(500 records) sets using stratified sampling to preserve the class ratio.
All random states were fixed at 42 for reproducibility.

% ────────────────────────────────────────────────────────────────────────────
\section{Evaluation and Model Comparison}
\label{sec:eval}

\subsection{Hold-Out Test Set Results}

Rather than using the default 0.5 decision threshold, each model's threshold
was tuned post-training under two strategies: (1) \textbf{F1-optimal} which
maximise the malicious-class F1 score, and (2) \textbf{Min-FN} which maximise
recall (drive false negatives to zero), then select the highest-precision
threshold among all that achieve that recall.
Tables~\ref{tab:test_f1} and~\ref{tab:test_minfn} report results for each
strategy on the held-out test set ($n=500$).

\begin{table}[H]
\centering
\caption{F1-optimal threshold --- hold-out test set (malicious class, $n=152$).}
\label{tab:test_f1}
\begin{tabular}{lrrrrr}
\toprule
\textbf{Model} & \textbf{Threshold} & \textbf{Accuracy} & \textbf{Precision} & \textbf{Recall} & \textbf{F1} \\
\midrule
Logistic Regression  & 0.384 & 0.9940 & 0.9869 & 0.9934 & 0.9902 \\
Random Forest        & 0.704 & 0.9940 & 1.0000 & 0.9803 & 0.9900 \\
\bottomrule
\end{tabular}
\end{table}

\begin{table}[H]
\centering
\caption{Min-FN threshold (zero false negatives) --- hold-out test set (malicious class, $n=152$).}
\label{tab:test_minfn}
\begin{tabular}{lrrrrr}
\toprule
\textbf{Model} & \textbf{Threshold} & \textbf{Accuracy} & \textbf{Precision} & \textbf{Recall} & \textbf{F1} \\
\midrule
Logistic Regression  & 0.192 & 0.9800 & 0.9383 & 1.0000 & 0.9682 \\
Random Forest        & 0.281 & 0.9920 & 0.9744 & 1.0000 & 0.9870 \\
\bottomrule
\end{tabular}
\end{table}

Both models achieve near-identical F1 scores ($\approx 0.990$) under the
F1-optimal strategy. Under the Min-FN strategy, both achieve perfect recall
(zero false negatives) at the cost of increased false positives: LR flags 10
benign flows as malicious versus 4 for RF. This makes RF the preferred choice
as it eliminates missed detections while
producing fewer false positive alerts than LR.

\subsection{ROC Curves}

The ROC curves (Appendix, Figure~\ref{fig:roc}) for both models are virtually
indistinguishable, hugging the top-left corner. This shows that they are able to correctly identify malicious traffic without being overly sensitive and causing false positives.

\subsection{Cross-Validation}

Table~\ref{tab:cv} reports 5-fold stratified cross-validation results.
Consistent with the hold-out evaluation, LR achieves the highest mean
accuracy ($99.56\% \pm 0.34\%$) and lowest variance, making it both the most
accurate and most stable classifier. Both models generalise well, with
standard deviations below 0.5\% for accuracy, indicating minimal overfitting.

\begin{table}[H]
\centering
\caption{5-fold stratified cross-validation summary (mean $\pm$ std).}
\label{tab:cv}
\begin{tabular}{lllll}
\toprule
\textbf{Model} & \textbf{Accuracy} & \textbf{Precision} & \textbf{Recall} & \textbf{F1} \\
\midrule
Logistic Regression
    & $0.9956\pm0.0034$ & $0.9935\pm0.0100$ & $0.9921\pm0.0064$ & $0.9928\pm0.0056$ \\
Random Forest
    & $0.9944\pm0.0039$ & $0.9935\pm0.0082$ & $0.9882\pm0.0113$ & $0.9908\pm0.0064$ \\
\bottomrule
\end{tabular}
\end{table}

\subsection{Feature Importance}

Random Forest feature importances (Appendix, Figure~\ref{fig:importance})
identify \texttt{packet\_count}, \texttt{num\_failed\_connections}, and
\texttt{failure\_rate} as the top three predictors by a wide margin (0.516,
0.239, and 0.098 respectively). Engineered features (\texttt{failure\_rate},
\texttt{bytes\_per\_packet}, \texttt{byte\_ratio}) appear among the top
predictors, validating the feature engineering step. The protocol indicators
(\texttt{protocol\_tcp}, \texttt{protocol\_udp}, \texttt{protocol\_icmp})
rank at the bottom, suggesting that attack traffic in this dataset is
protocol-agnostic and relies more on behavioural anomalies than
protocol-specific exploitation.


% ────────────────────────────────────────────────────────────────────────────
\section{Forensic Implications and Ethical Concerns}
\label{sec:forensics}

\subsection{False Negatives in a Forensic Context}

A false negative, a malicious flow classified as benign, is the more
consequential error in digital forensics. Missed detections may
allow attackers to persist undetected, potentially obscuring the full scope
of a breach. Under the Min-FN
threshold strategy, both models achieve zero false negatives on the test set,
though this likely will not be the case in real-world scenarios.

\subsection{False Positives and Alert Fatigue}

False positives have the very real potential to cause analyst fatigue. If a security department is
handling millions of flows per day, even a 1\% false-positive rate would
generate thousands of alerts per hour. Under the F1-optimal threshold, the models produce 0--2 false positives per
348 benign test records ($\approx 0\%$--$0.6\%$), which may be acceptable for
a first-pass automated tool. Under the Min-FN threshold the rate
rises to 4--10 false positives ($\approx 1.1\%$--$2.9\%$), which may be a worthwhile
trade-off provided flagged flows are reviewed by a human analyst before any
action is taken.

\subsection{Ethical and Legal Concerns}

\textbf{Privacy:} Network flow analysis reveals communication patterns between
individuals and organisations. Aggregate flow metadata may expose
sensitive behavioural information (medical service connections,
political affiliations, etc.). Users of a model like this should ensure that data collection
and storage comply with local privacy laws. 

\textbf{Bias and Fairness:} Training data labelled by a malicious traffic detection model may systematically mis-label traffic from underrepresented network
configurations or geographic regions. Models trained on such data may
present bias, leading to a disproportionate scrutiny of
certain users.


% ────────────────────────────────────────────────────────────────────────────
\section{Conclusion}

This project demonstrated that supervised machine learning applied to
engineered network flow features can classify malicious traffic with
exceptional accuracy. Both models achieved near-identical F1 scores of
$\approx 0.990$ under the F1-optimal threshold, with Random Forest preferred
in a forensic context due to its ability to eliminate false negatives while
producing fewer false positives than Logistic Regression. Both models
achieved near-identical performance in 5-fold cross-validation, confirming
stable generalisation. The forensic analysis highlights that while the
models are highly effective as automated triage tools, human oversight
remains indispensable to contextualise predictions, manage false-positive
burden, and ensure legally defensible conclusions.

% ────────────────────────────────────────────────────────────────────────────
\section{Code}
\label{sec:code}

All Python source code for data loading, feature engineering, model training,
evaluation, and figure generation is available in the project repository:\\
\url{https://github.com/jhelani-stokes/CSCI-4540-Final}

% ────────────────────────────────────────────────────────────────────────────
\clearpage
\appendix
\section{Visualizations and Evaluation Metrics}

\begin{figure}[H]
    \centering
    \hfill
    \begin{subfigure}[b]{0.85\textwidth}
        \includegraphics[width=\textwidth]{figures/correlation_heatmap.png}
        \caption{Pearson correlation heatmap.}
        \label{fig:heatmap}
    \end{subfigure}
    \caption{Dataset overview.}
\end{figure}

\begin{figure}[H]
    \centering
    \begin{subfigure}[b]{0.48\textwidth}
        \includegraphics[width=\textwidth]{figures/cm_logistic_regression_f1.png}
        \caption{LR --- F1-optimal threshold.}
        \label{fig:cm_lr_f1}
    \end{subfigure}
    \hfill
    \begin{subfigure}[b]{0.48\textwidth}
        \includegraphics[width=\textwidth]{figures/cm_logistic_regression_minfn.png}
        \caption{LR --- Min-FN threshold.}
        \label{fig:cm_lr_minfn}
    \end{subfigure}
    \vspace{0.5em}
    \begin{subfigure}[b]{0.48\textwidth}
        \includegraphics[width=\textwidth]{figures/cm_random_forest_f1.png}
        \caption{RF --- F1-optimal threshold.}
        \label{fig:cm_rf_f1}
    \end{subfigure}
    \hfill
    \begin{subfigure}[b]{0.48\textwidth}
        \includegraphics[width=\textwidth]{figures/cm_random_forest_minfn.png}
        \caption{RF --- Min-FN threshold.}
        \label{fig:cm_rf_minfn}
    \end{subfigure}
    \caption{Confusion matrices on the hold-out test set under both threshold strategies.}
    \label{fig:cms}
\end{figure}

\begin{figure}[H]
    \centering
    \includegraphics[width=0.7\textwidth]{figures/roc_curves.png}
    \caption{ROC curves.}
    \label{fig:roc}
\end{figure}

\begin{figure}[H]
    \centering
    \includegraphics[width=0.85\textwidth]{figures/cv_comparison.png}
    \caption{5-fold cross-validation comparison (mean $\pm$ std).}
    \label{fig:cv}
\end{figure}

\begin{figure}[H]
    \centering
    \includegraphics[width=0.85\textwidth]{figures/feature_importance.png}
    \caption{Top 15 feature importances from the Random Forest classifier.}
    \label{fig:importance}
\end{figure}

\end{document}
